\documentclass[11pt]{article}

\usepackage{sectsty}
\usepackage{graphicx}
\usepackage{amsmath}

% Margins
\topmargin=-0.45in
\evensidemargin=0in
\oddsidemargin=0in
\textwidth=6.5in
\textheight=9.0in
\headsep=0.25in

\title{Generalized note on HMM with AR(1)}
% \author{ Author }
% \date{\today}

\begin{document}
\maketitle	
\pagebreak

% Optional TOC
% \tableofcontents
% \pagebreak

%--Paper--

\section{Introduction}

Let $s_t\in \{1, 2, ..., K\}$ denote the hidden state at time $t$. The states 
follow a first-order Markov chain with transition probabilites 
$$
P(s_t=j\vert s_{t-1}=i) = p_{ij},
$$
with transition matrix
$$
\begin{pmatrix}
    p_{11} & \cdots & p_{1K}\\
    \vdots & \ddots & \vdots \\
    p_{K1} & \cdots & p_{KK}
\end{pmatrix},
$$
where each row sum must sum to one:
$$
\sum_{j=1}^K p_{ij} = 1.
$$
The initial distribution is $\pi=(\pi_1, \ldots, \pi_K)$ with $P(s_1=i)=\pi_i$
and $\sum_{i=1}^K \pi_i = 1$. Then, conditional on the state $s_t$, the observations
follow a state-dependent autoregressive model, i.e.
$$
y_t = \beta_{s_t}y_{t-1} + \epsilon_t, 
$$
where $\epsilon_t \sim N(0, \sigma^2_{s_t})$. Thus the conditional distribution of
the observation is 
$$
y_t \vert (s_t=i, y_{t-1}) \sim N(\beta_i y_{t-1}, \sigma^2_i),
$$

and the corresponding density is 
$$
f(y_t \vert s_t = i, y_{t-1}) = \frac{1}{\sqrt{2\pi \sigma^2_i}} 
\operatorname{exp}\left(- \frac{(y_t - \beta_i y_{t-1})^2}{2\sigma^2_i}\right)
$$

\section{Likelihood}
Now, if all states $s_{1:T}$ were observed, the likelihood would be
$$
L = p(s_1)\prod_{t=2}^T p(s_t \mid s_{t-1}) \prod_{t=2}^T f(y_t \mid s_t, y_{t-1}),
$$
and the complete-data log-likelihood is therefore
$$
\log L = \log p(s_1) + \sum_{t=2}^T \log p(s_t \mid s_{t-1}) + \sum_{t=2}^T \log f(y_t \mid s_t, y_{t-1}),
$$
with $f(y_t \mid s_t = i, y_{t-1})$ as defined above. Let $y_{1:T} = (y_1, \ldots, y_T)$ denote the observed time series and $s_t \in \{1, \ldots, K\}$ the hidden states. Since the states are unobserved, the likelihood of the data is obtained by summing over all possible state sequences:
$$
L = P(y_{1:T}) = \sum_{s_1=1}^K \sum_{s_2=1}^K \cdots \sum_{s_T=1}^K P(y_{1:T}, s_{1:T}),
$$
where
$$
P(y_{1:T}, s_{1:T}) = p(s_1)\prod_{t=2}^T p(s_t \mid s_{t-1}) \prod_{t=2}^T f(y_t \mid s_t, y_{t-1}),
$$
so the observed likelihood is
$$
L = \sum_{s_1=1}^K \cdots \sum_{s_T=1}^K p(s_1)\prod_{t=2}^T p(s_t \mid s_{t-1}) \prod_{t=2}^T f(y_t \mid s_t, y_{t-1}).
$$

With $K$ states this involves $K^T$ terms, which becomes computationally infeasible for large $T$.

\section{Forward Algorithm}
To compute the likelihood efficiently we introduce the forward probabilities
$$
\alpha_t(i) = P(y_{1:t}, s_t = i).
$$
This quantity represents the joint probability of observing the data up to time $t$ and being in state $i$ at time $t$. For the AR(1) observation model, we condition on the first observation $y_1$ and treat it as fixed (or assume a known $y_0$), so we initialise with
$$
\alpha_1(i) = \pi_i f(y_1 \mid s_1 = i, y_0).
$$
Then, by using the law of total probability, we have
$$
\alpha_t(j) = P(y_{1:t}, s_t = j) = \sum_{i=1}^K P(y_{1:t}, s_t = j, s_{t-1} = i),
$$
and by factorising the joint probability,
$$
\alpha_t(j) = \sum_{i=1}^K P(y_t \mid s_t = j, y_{t-1}) P(s_t = j \mid s_{t-1} = i) P(y_{1:t-1}, s_{t-1} = i),
$$
which gives the recursion
$$
\alpha_t(j) =\left(\sum_{i=1}^K \alpha_{t-1}(i) p_{ij} \right)\, f(y_t \mid s_t = j, y_{t-1}), \quad j \in \{1, \ldots, K\}.
$$
The likelihood of the observed data is obtained by marginalising over the final state:
$$
L = P(y_{1:T}) = \sum_{i=1}^K \alpha_T(i).
$$
In practice, the forward probabilities $\alpha_t(i)$ shrink exponentially in $T$, leading to numerical underflow. A standard trick is to work with the log-likelihood directly by rescaling at each step, i.e. define
$$
c_t = \sum_{i=1}^K \alpha_t(i), \quad \tilde{\alpha}_t(i) = \frac{\alpha_t(i)}{c_t},
$$
then
$$
\log L = \sum_{t=1}^T \log c_t.
$$
Now, the unknown parameters $(\beta_1, \ldots, \beta_K, \sigma_1, \ldots, \sigma_K, P, \pi)$ 
can be found by maximising $\log L$ using numerical optimisation. 
Note that it is advisable to use unconstrained parameterisations: 
$\sigma_i = \exp(\eta_i)$ ensures $\sigma_i > 0$; transition probabilities 
$p_{ij} = \exp(a_{ij}) / \sum_{l=1}^K \exp(a_{il})$ ensure $p_{ij} \in (0,1)$ 
and that each row sums to one; and $\beta_i = (1 - \exp(-b_i))/(1 + \exp(-b_i))$ 
ensures $\beta_i \in (-1,1)$, which corresponds to stationarity of the AR(1) 
process in each state.

\section{Filtered State Probabilites}
Suppose we observe a sequence of data $y_1, y_2, \ldots$ sequentially. We are interested in the filtered state probabilities
$$
P(s_t = i \mid y_{1:t}),
$$
i.e. the probability of being in state $i$ given all observations up to time $t$. To estimate this, recall that $\alpha_t(i) = P(y_{1:t}, s_t = i)$. Initialise $\alpha_1(i) = \pi_i f(y_1 \mid s_1 = i, y_0)$, and then by recursion for $t \geq 2$:
$$
\alpha_t(j) = \sum_{i=1}^K \alpha_{t-1}(i) p_{ij} \, f(y_t \mid s_t = j, y_{t-1}), \quad j \in \{1, \ldots, K\}.
$$
The filtered state probabilities are obtained by normalisation:
$$
P(s_t = i \mid y_{1:t}) = \frac{\alpha_t(i)}{\sum_{j=1}^K \alpha_t(j)}.
$$

\section{ARIMAX, ARCH, GARCH}

\begin{center}
\textbf{Warning:} This section is generated by AI.
\end{center}

The hidden Markov model framework presented above can be extended naturally to 
a wide class of time series models by allowing the model parameters to depend 
on the hidden state. The key idea is that, conditional on the state $s_t = i$, 
the observations follow a state-dependent model with parameters specific to regime $i$.

\subsection{ARIMAX models}

An immediate extension is to autoregressive models with exogenous variables. 
In this case, conditional on $s_t = i$, the observations follow an ARIMAX-type model:
$$
y_t = c_i + \sum_{l=1}^p \phi_{l,i} y_{t-l} + \sum_{m=1}^q \theta_{m,i} 
\varepsilon_{t-m} + \beta_i^\top x_t + \varepsilon_t,
$$
where $\varepsilon_t \sim N(0, \sigma_i^2)$ and $x_t$ denotes a vector of 
exogenous variables. Each regime is thus characterised by its own set of 
autoregressive, moving average and regression coefficients, as well as its own variance.

\subsection{ARCH and GARCH models}

The framework can also be extended to models with time-varying volatility. 
In a regime-switching ARCH model, the conditional variance depends on both past 
observations and the current state:
$$
y_t = \sigma_t \varepsilon_t, \quad \varepsilon_t \sim N(0,1),
$$
$$
\sigma_t^2 = \alpha_{0,i} + \alpha_{1,i} y_{t-1}^2, \quad \text{if } s_t = i.
$$

Similarly, in a regime-switching GARCH model, the conditional variance follows 
a GARCH-type recursion within each state:
$$
\sigma_t^2 = \alpha_{0,i} + \alpha_{1,i} y_{t-1}^2 + \beta_{1,i} \sigma_{t-1}^2, 
\quad \text{if } s_t = i.
$$

\subsection{General formulation}

More generally, the hidden Markov model framework can be written as
$$
y_t \mid (s_t = i, \mathcal{F}_{t-1}) \sim f(\cdot; \theta_i, \mathcal{F}_{t-1}),
$$
where $\mathcal{F}_{t-1}$ denotes the information set up to time $t-1$, and 
$\theta_i$ is a set of state-dependent parameters. The choice of $f$ determines 
the specific model, such as AR(1), ARIMAX, ARCH or GARCH.

This formulation highlights that the hidden Markov model provides a flexible 
way to capture structural changes in time series by allowing different regimes 
to follow different dynamic models.


\section{Order Estimation in Hidden Markov Models}

\begin{center}
\textbf{Warning:} This section is generated by AI.
\end{center}

A key modelling choice in Hidden Markov Models (HMMs) is the number of hidden states, denoted by $K$. This is often referred to as the \emph{order estimation problem}. Selecting an appropriate value of $K$ is crucial, as too few states may lead to underfitting, while too many states may result in overfitting and poor predictive performance.

\subsection{Likelihood-Based Model Selection}

Let $\{y_t\}_{t=1}^n$ denote the observed time series. For a given number of states $K$, the HMM is estimated by maximizing the log-likelihood
\[
\ell_K = \log p(y_1, \dots, y_n \mid \hat{\theta}_K),
\]
where $\hat{\theta}_K$ denotes the estimated parameters of the model with $K$ states.

The likelihood can be efficiently computed using the forward algorithm. For each candidate model (i.e., each value of $K$), we obtain the corresponding maximized log-likelihood $\ell_K$.

\subsection{Information Criteria}

To balance model fit and complexity, information criteria such as the Akaike Information Criterion (AIC) and Bayesian Information Criterion (BIC) are commonly used.

For a model with $k_K$ free parameters and sample size $n$, these are defined as
\[
\text{AIC}(K) = -2 \ell_K + 2 k_K,
\]
\[
\text{BIC}(K) = -2 \ell_K + k_K \log n.
\]

The preferred model is the one that minimizes the chosen criterion.

\subsection{Number of Parameters in an HMM}

The number of free parameters $k_K$ in an HMM depends on the model specification. For a Gaussian HMM with $K$ states, we typically have:

\begin{itemize}
    \item Initial state probabilities: $K - 1$
    \item Transition matrix: $K(K - 1)$
    \item Emission parameters: depends on the model
\end{itemize}

For example, in an AR(1)-HMM with state-dependent parameters $(\mu_i, \phi_i, \sigma_i^2)$ for each state $i = 1, \dots, K$, the number of emission parameters is $3K$.

Thus, the total number of parameters is
\[
k_K = (K - 1) + K(K - 1) + 3K.
\]

\subsection{Practical Procedure}

In practice, order estimation is performed as follows:

\begin{enumerate}
    \item Choose a range of candidate values for $K$ (e.g., $K = 1, 2, 3, 4, 5$).
    \item Estimate the HMM for each $K$ using maximum likelihood.
    \item Compute $\ell_K$, AIC$(K)$, and BIC$(K)$.
    \item Select the model with the lowest BIC (or AIC).
\end{enumerate}

BIC is often preferred in this context due to its stronger penalty for model complexity, which helps avoid overfitting.

\subsection{Remarks}

It is important to note that standard likelihood ratio tests are generally not valid for determining the number of states in HMMs, due to non-standard asymptotic behavior and identifiability issues. As a result, information criteria and Bayesian methods are typically preferred in practice.

Finally, model selection should not rely solely on information criteria. It is also important to assess whether the inferred states are interpretable and whether the model provides improved predictive performance.
\pagebreak
% \section{Section 2}
% Lorem Ipsum \\

%--/Paper--

\end{document}